% --- Archon physics formalization source begin ---
% archon:physics
% archon:covers PhyXMiniProblems/problem_phyx_mini_0718.lean
% archon:source-report reports/phyx_mini/problem_phyx_mini_0718.source.json

\chapter{Physics problem phyx\_mini\_0718}
\label{ch:PhyXMiniProblems_problem_phyx_mini_0718}

\paragraph{Problem source.}
\#\# Physical scenario

An \textbackslash{}(83\textbackslash{},\textbackslash{}mathrm\{kg\}\textbackslash{}) student hangs from a bungee cord with spring constant \textbackslash{}(270\textbackslash{},\textbackslash{}mathrm\{N/m\}\textbackslash{}). The student is pulled down to a point where the cord is \textbackslash{}(5.0\textbackslash{},\textbackslash{}mathrm\{m\}\textbackslash{}) longer than its unstretched length, then released.

\#\# Question

What is his velocity \textbackslash{}(2.0\textbackslash{},\textbackslash{}mathrm\{s\}\textbackslash{}) later?

\#\# Answer choices

- A: A: -2.0m/s

- B: B: -1.8m/s

- C: C: -1.6m/s

- D: D: -1.4m/s

\#\# Figure caption (auxiliary; use the image as primary evidence)

The image illustrates a bungee jumping scenario where a bungee cord is modeled as a spring. The diagram includes the following elements:

1. **Bungee Cord**: 
   - Represented as a coiled spring attached to a ceiling.

2. **Spring Constant**: 
   - Labeled as "270 N/m" next to the spring on the left.

3. **Distances**:
   - Two horizontal dashed lines labeled "Equilibrium" and "Release," with a vertical arrow between them labeled "5.0 m."

4. **Person**:
   - A stick figure at the end of the bungee cord on the right side, labeled "83 kg."

5. **Oscillation Diagram**:
   - To the right, a vertical line represents oscillation with markers labeled "A," "0," and "-A," showing upward and downward arrows with the text "Oscillation."

6. **Text**:
   - A note states, "The bungee cord is modeled as a spring."

This diagram visually explains the concept of modeling a bungee jump as a spring-mass system, highlighting key parameters like spring constant, mass, and oscillation points.

\#\# Dataset metadata

Dynamics; Physical Model Grounding Reasoning, Predictive Reasoning

\paragraph{Recorded answer/context.}
C: C: -1.6m/s

\paragraph{Figure/image path.}
/root/proposal\_for\_physic/science-mango/phyx\_mini\_run/phyx\_data/test\_image/718.png

\paragraph{Formalization target.}
create a compiling Lean file with sorry bodies at `PhyXMiniProblems/problem\_phyx\_mini\_0718.lean`. The Lean declarations must preserve the physical quantities, dimensions or dimensional roles, figure labels, governing-law hypotheses, and final relation expressed by this problem.
Use Mathlib/Physlib names found through LeanExplore where available. If a physics API is missing, introduce faithful local abstractions rather than scalar placeholder aliases.

\begin{theorem}[Physics formalization target]
\label{thm:physics:phyx_mini_0718:target}
\lean{PhyXMiniProblems.ProblemPhyXMini0718.bungee_velocity_after_two_seconds_matches_answer_C}
\uses{def:physics:phyx-mini-0718:phyxminiproblems-problemphyxmini0718-velocityinmeterspersecond, def:physics:phyx-mini-0718:phyxminiproblems-problemphyxmini0718-bungeeoscillationsetup, def:physics:phyx-mini-0718:phyxminiproblems-problemphyxmini0718-matchesproblemdata, def:physics:phyx-mini-0718:phyxminiproblems-problemphyxmini0718-matchessourcefigure, def:physics:phyx-mini-0718:phyxminiproblems-problemphyxmini0718-satisfiesidealbungeephysics, def:physics:phyx-mini-0718:phyxminiproblems-problemphyxmini0718-answerchoice, def:physics:phyx-mini-0718:phyxminiproblems-problemphyxmini0718-answerchoice-velocitymeterspersecond}
The assigned autoformalize agent should translate this physics problem into Lean declarations in the covered file, with theorem and lemma proof bodies written as `by sorry`.
\end{theorem}
\begin{proof}
This is an autoformalization task, not a proof task. Produce statements that can later be proved by the physics prover without weakening the physical model.
\end{proof}

% --- Archon declaration topology begin ---
\section{Declaration topology}
\begin{definition}[velocity Dimension]
  \label{def:physics:phyx-mini-0718:phyxminiproblems-problemphyxmini0718-velocitydimension}
  \lean{PhyXMiniProblems.ProblemPhyXMini0718.velocityDimension}
  The physical dimension of velocity, L T⁻¹
\end{definition}
\begin{proof}
  This is the declared model definition of velocity Dimension.
\end{proof}
\begin{definition}[acceleration Dimension]
  \label{def:physics:phyx-mini-0718:phyxminiproblems-problemphyxmini0718-accelerationdimension}
  \lean{PhyXMiniProblems.ProblemPhyXMini0718.accelerationDimension}
  The physical dimension of acceleration, L T⁻²
\end{definition}
\begin{proof}
  This is the declared model definition of acceleration Dimension.
\end{proof}
\begin{definition}[spring Stiffness Dimension]
  \label{def:physics:phyx-mini-0718:phyxminiproblems-problemphyxmini0718-springstiffnessdimension}
  \lean{PhyXMiniProblems.ProblemPhyXMini0718.springStiffnessDimension}
  The physical dimension of linear spring stiffness, M T⁻²
\end{definition}
\begin{proof}
  This is the declared model definition of spring Stiffness Dimension.
\end{proof}
\begin{definition}[Mass Quantity]
  \label{def:physics:phyx-mini-0718:phyxminiproblems-problemphyxmini0718-massquantity}
  \lean{PhyXMiniProblems.ProblemPhyXMini0718.MassQuantity}
  A nonnegative physical mass
\end{definition}
\begin{proof}
  This is the declared model definition of Mass Quantity.
\end{proof}
\begin{definition}[Length Quantity]
  \label{def:physics:phyx-mini-0718:phyxminiproblems-problemphyxmini0718-lengthquantity}
  \lean{PhyXMiniProblems.ProblemPhyXMini0718.LengthQuantity}
  A nonnegative physical length
\end{definition}
\begin{proof}
  This is the declared model definition of Length Quantity.
\end{proof}
\begin{definition}[Time Quantity]
  \label{def:physics:phyx-mini-0718:phyxminiproblems-problemphyxmini0718-timequantity}
  \lean{PhyXMiniProblems.ProblemPhyXMini0718.TimeQuantity}
  A nonnegative physical duration elapsed since release
\end{definition}
\begin{proof}
  This is the declared model definition of Time Quantity.
\end{proof}
\begin{definition}[Spring Stiffness Quantity]
  \label{def:physics:phyx-mini-0718:phyxminiproblems-problemphyxmini0718-springstiffnessquantity}
  \lean{PhyXMiniProblems.ProblemPhyXMini0718.SpringStiffnessQuantity}
  \uses{def:physics:phyx-mini-0718:phyxminiproblems-problemphyxmini0718-springstiffnessdimension}
  A nonnegative physical linear-spring stiffness
\end{definition}
\begin{proof}
  This is the declared model definition of Spring Stiffness Quantity.
\end{proof}
\begin{definition}[Acceleration Magnitude Quantity]
  \label{def:physics:phyx-mini-0718:phyxminiproblems-problemphyxmini0718-accelerationmagnitudequantity}
  \lean{PhyXMiniProblems.ProblemPhyXMini0718.AccelerationMagnitudeQuantity}
  \uses{def:physics:phyx-mini-0718:phyxminiproblems-problemphyxmini0718-accelerationdimension}
  A nonnegative physical acceleration magnitude
\end{definition}
\begin{proof}
  This is the declared model definition of Acceleration Magnitude Quantity.
\end{proof}
\begin{definition}[Signed Length Quantity]
  \label{def:physics:phyx-mini-0718:phyxminiproblems-problemphyxmini0718-signedlengthquantity}
  \lean{PhyXMiniProblems.ProblemPhyXMini0718.SignedLengthQuantity}
  A signed vertical displacement from equilibrium
\end{definition}
\begin{proof}
  This is the declared model definition of Signed Length Quantity.
\end{proof}
\begin{definition}[Signed Velocity Quantity]
  \label{def:physics:phyx-mini-0718:phyxminiproblems-problemphyxmini0718-signedvelocityquantity}
  \lean{PhyXMiniProblems.ProblemPhyXMini0718.SignedVelocityQuantity}
  \uses{def:physics:phyx-mini-0718:phyxminiproblems-problemphyxmini0718-velocitydimension}
  A signed vertical velocity, positive in the displayed upward direction
\end{definition}
\begin{proof}
  This is the declared model definition of Signed Velocity Quantity.
\end{proof}
\begin{definition}[nonnegative SIReadout]
  \label{def:physics:phyx-mini-0718:phyxminiproblems-problemphyxmini0718-nonnegativesireadout}
  \lean{PhyXMiniProblems.ProblemPhyXMini0718.nonnegativeSIReadout}
  Read a nonnegative dimensionful quantity in coherent SI units
\end{definition}
\begin{proof}
  This is the declared model definition of nonnegative SIReadout.
\end{proof}
\begin{definition}[signed SIReadout]
  \label{def:physics:phyx-mini-0718:phyxminiproblems-problemphyxmini0718-signedsireadout}
  \lean{PhyXMiniProblems.ProblemPhyXMini0718.signedSIReadout}
  Read a signed dimensionful quantity in coherent SI units
\end{definition}
\begin{proof}
  This is the declared model definition of signed SIReadout.
\end{proof}
\begin{definition}[mass In Kilograms]
  \label{def:physics:phyx-mini-0718:phyxminiproblems-problemphyxmini0718-massinkilograms}
  \lean{PhyXMiniProblems.ProblemPhyXMini0718.massInKilograms}
  \uses{def:physics:phyx-mini-0718:phyxminiproblems-problemphyxmini0718-massquantity, def:physics:phyx-mini-0718:phyxminiproblems-problemphyxmini0718-nonnegativesireadout}
  Kilogram readout of the student's physical mass
\end{definition}
\begin{proof}
  This is the declared model definition of mass In Kilograms.
\end{proof}
\begin{definition}[length In Meters]
  \label{def:physics:phyx-mini-0718:phyxminiproblems-problemphyxmini0718-lengthinmeters}
  \lean{PhyXMiniProblems.ProblemPhyXMini0718.lengthInMeters}
  \uses{def:physics:phyx-mini-0718:phyxminiproblems-problemphyxmini0718-lengthquantity, def:physics:phyx-mini-0718:phyxminiproblems-problemphyxmini0718-nonnegativesireadout}
  Metre readout of a nonnegative physical length
\end{definition}
\begin{proof}
  This is the declared model definition of length In Meters.
\end{proof}
\begin{definition}[time In Seconds]
  \label{def:physics:phyx-mini-0718:phyxminiproblems-problemphyxmini0718-timeinseconds}
  \lean{PhyXMiniProblems.ProblemPhyXMini0718.timeInSeconds}
  \uses{def:physics:phyx-mini-0718:phyxminiproblems-problemphyxmini0718-timequantity, def:physics:phyx-mini-0718:phyxminiproblems-problemphyxmini0718-nonnegativesireadout}
  Second readout of a physical elapsed duration
\end{definition}
\begin{proof}
  This is the declared model definition of time In Seconds.
\end{proof}
\begin{definition}[stiffness In Newtons Per Meter]
  \label{def:physics:phyx-mini-0718:phyxminiproblems-problemphyxmini0718-stiffnessinnewtonspermeter}
  \lean{PhyXMiniProblems.ProblemPhyXMini0718.stiffnessInNewtonsPerMeter}
  \uses{def:physics:phyx-mini-0718:phyxminiproblems-problemphyxmini0718-springstiffnessquantity, def:physics:phyx-mini-0718:phyxminiproblems-problemphyxmini0718-nonnegativesireadout}
  Newton-per-metre readout of the cord's spring stiffness
\end{definition}
\begin{proof}
  This is the declared model definition of stiffness In Newtons Per Meter.
\end{proof}
\begin{definition}[acceleration In Meters Per Second Squared]
  \label{def:physics:phyx-mini-0718:phyxminiproblems-problemphyxmini0718-accelerationinmeterspersecondsquared}
  \lean{PhyXMiniProblems.ProblemPhyXMini0718.accelerationInMetersPerSecondSquared}
  \uses{def:physics:phyx-mini-0718:phyxminiproblems-problemphyxmini0718-accelerationmagnitudequantity, def:physics:phyx-mini-0718:phyxminiproblems-problemphyxmini0718-nonnegativesireadout}
  Metre-per-second-squared readout of an acceleration magnitude
\end{definition}
\begin{proof}
  This is the declared model definition of acceleration In Meters Per Second Squared.
\end{proof}
\begin{definition}[signed Length In Meters]
  \label{def:physics:phyx-mini-0718:phyxminiproblems-problemphyxmini0718-signedlengthinmeters}
  \lean{PhyXMiniProblems.ProblemPhyXMini0718.signedLengthInMeters}
  \uses{def:physics:phyx-mini-0718:phyxminiproblems-problemphyxmini0718-signedlengthquantity, def:physics:phyx-mini-0718:phyxminiproblems-problemphyxmini0718-signedsireadout}
  Metre readout of a signed upward-positive displacement
\end{definition}
\begin{proof}
  This is the declared model definition of signed Length In Meters.
\end{proof}
\begin{definition}[velocity In Meters Per Second]
  \label{def:physics:phyx-mini-0718:phyxminiproblems-problemphyxmini0718-velocityinmeterspersecond}
  \lean{PhyXMiniProblems.ProblemPhyXMini0718.velocityInMetersPerSecond}
  \uses{def:physics:phyx-mini-0718:phyxminiproblems-problemphyxmini0718-signedvelocityquantity, def:physics:phyx-mini-0718:phyxminiproblems-problemphyxmini0718-signedsireadout}
  Metre-per-second readout of a signed upward-positive velocity
\end{definition}
\begin{proof}
  This is the declared model definition of velocity In Meters Per Second.
\end{proof}
\begin{definition}[Figure Object]
  \label{def:physics:phyx-mini-0718:phyxminiproblems-problemphyxmini0718-figureobject}
  \lean{PhyXMiniProblems.ProblemPhyXMini0718.FigureObject}
  Physical and graphical objects visibly present in the supplied image
\end{definition}
\begin{proof}
  This is the declared model definition of Figure Object.
\end{proof}
\begin{definition}[Figure Text Label]
  \label{def:physics:phyx-mini-0718:phyxminiproblems-problemphyxmini0718-figuretextlabel}
  \lean{PhyXMiniProblems.ProblemPhyXMini0718.FigureTextLabel}
  Text and mathematical labels printed in the supplied image
\end{definition}
\begin{proof}
  This is the declared model definition of Figure Text Label.
\end{proof}
\begin{definition}[Figure Level]
  \label{def:physics:phyx-mini-0718:phyxminiproblems-problemphyxmini0718-figurelevel}
  \lean{PhyXMiniProblems.ProblemPhyXMini0718.FigureLevel}
  Distinguished vertical levels in the supplied bungee drawing
\end{definition}
\begin{proof}
  This is the declared model definition of Figure Level.
\end{proof}
\begin{definition}[Bungee Figure Readout]
  \label{def:physics:phyx-mini-0718:phyxminiproblems-problemphyxmini0718-bungeefigurereadout}
  \lean{PhyXMiniProblems.ProblemPhyXMini0718.BungeeFigureReadout}
  \uses{def:physics:phyx-mini-0718:phyxminiproblems-problemphyxmini0718-figureobject, def:physics:phyx-mini-0718:phyxminiproblems-problemphyxmini0718-figuretextlabel, def:physics:phyx-mini-0718:phyxminiproblems-problemphyxmini0718-figurelevel}
  The categorical incidences and scalar labels transcribed from the primary image. verticalCoordinate is only a schematic figure coordinate; it is not a replacement for any dimensionful length in the physical setup
\end{definition}
\begin{proof}
  This is the declared model definition of Bungee Figure Readout.
\end{proof}
\begin{definition}[Vertical Coordinate Convention]
  \label{def:physics:phyx-mini-0718:phyxminiproblems-problemphyxmini0718-verticalcoordinateconvention}
  \lean{PhyXMiniProblems.ProblemPhyXMini0718.VerticalCoordinateConvention}
  The sign and origin convention shown by the vertical y axis
\end{definition}
\begin{proof}
  This is the declared model definition of Vertical Coordinate Convention.
\end{proof}
\begin{definition}[Cord Model]
  \label{def:physics:phyx-mini-0718:phyxminiproblems-problemphyxmini0718-cordmodel}
  \lean{PhyXMiniProblems.ProblemPhyXMini0718.CordModel}
  Idealization used for the bungee cord over the modeled interval
\end{definition}
\begin{proof}
  This is the declared model definition of Cord Model.
\end{proof}
\begin{definition}[Motion Regime]
  \label{def:physics:phyx-mini-0718:phyxminiproblems-problemphyxmini0718-motionregime}
  \lean{PhyXMiniProblems.ProblemPhyXMini0718.MotionRegime}
  Motion regime assumed after the student is released
\end{definition}
\begin{proof}
  This is the declared model definition of Motion Regime.
\end{proof}
\begin{definition}[Bungee Oscillation Setup]
  \label{def:physics:phyx-mini-0718:phyxminiproblems-problemphyxmini0718-bungeeoscillationsetup}
  \lean{PhyXMiniProblems.ProblemPhyXMini0718.BungeeOscillationSetup}
  \uses{def:physics:phyx-mini-0718:phyxminiproblems-problemphyxmini0718-massquantity, def:physics:phyx-mini-0718:phyxminiproblems-problemphyxmini0718-lengthquantity, def:physics:phyx-mini-0718:phyxminiproblems-problemphyxmini0718-timequantity, def:physics:phyx-mini-0718:phyxminiproblems-problemphyxmini0718-springstiffnessquantity, def:physics:phyx-mini-0718:phyxminiproblems-problemphyxmini0718-accelerationmagnitudequantity, def:physics:phyx-mini-0718:phyxminiproblems-problemphyxmini0718-signedlengthquantity, def:physics:phyx-mini-0718:phyxminiproblems-problemphyxmini0718-signedvelocityquantity, def:physics:phyx-mini-0718:phyxminiproblems-problemphyxmini0718-bungeefigurereadout, def:physics:phyx-mini-0718:phyxminiproblems-problemphyxmini0718-verticalcoordinateconvention, def:physics:phyx-mini-0718:phyxminiproblems-problemphyxmini0718-cordmodel, def:physics:phyx-mini-0718:phyxminiproblems-problemphyxmini0718-motionregime}
  The physical quantities and response functions for the bungee experiment. unstretchedToReleaseExtension is the full 5.0 m cord extension visible in the image. It is deliberately distinct from oscillationAmplitude, which is the release-to-equilibrium distance after the gravitational equilibrium shift. The requested velocity is not stored as a special numerical field
\end{definition}
\begin{proof}
  This is the declared model definition of Bungee Oscillation Setup.
\end{proof}
\begin{definition}[Matches Problem Data]
  \label{def:physics:phyx-mini-0718:phyxminiproblems-problemphyxmini0718-matchesproblemdata}
  \lean{PhyXMiniProblems.ProblemPhyXMini0718.MatchesProblemData}
  \uses{def:physics:phyx-mini-0718:phyxminiproblems-problemphyxmini0718-massinkilograms, def:physics:phyx-mini-0718:phyxminiproblems-problemphyxmini0718-lengthinmeters, def:physics:phyx-mini-0718:phyxminiproblems-problemphyxmini0718-timeinseconds, def:physics:phyx-mini-0718:phyxminiproblems-problemphyxmini0718-stiffnessinnewtonspermeter, def:physics:phyx-mini-0718:phyxminiproblems-problemphyxmini0718-accelerationinmeterspersecondsquared, def:physics:phyx-mini-0718:phyxminiproblems-problemphyxmini0718-bungeeoscillationsetup}
  Numerical values given in the prose, together with the standard near-Earth gravity calibration used by the recorded multiple-choice calculation. No velocity value or answer-choice assertion occurs here
\end{definition}
\begin{proof}
  This is the declared model definition of Matches Problem Data.
\end{proof}
\begin{definition}[Matches Source Figure]
  \label{def:physics:phyx-mini-0718:phyxminiproblems-problemphyxmini0718-matchessourcefigure}
  \lean{PhyXMiniProblems.ProblemPhyXMini0718.MatchesSourceFigure}
  \uses{def:physics:phyx-mini-0718:phyxminiproblems-problemphyxmini0718-massinkilograms, def:physics:phyx-mini-0718:phyxminiproblems-problemphyxmini0718-lengthinmeters, def:physics:phyx-mini-0718:phyxminiproblems-problemphyxmini0718-stiffnessinnewtonspermeter, def:physics:phyx-mini-0718:phyxminiproblems-problemphyxmini0718-bungeeoscillationsetup}
  Transcription of the source image and coherence of its three numerical labels with the corresponding dimensionful quantities. The top endpoint of the 5.0 m arrow is the unstretched-cord level, not the equilibrium line
\end{definition}
\begin{proof}
  This is the declared model definition of Matches Source Figure.
\end{proof}
\begin{definition}[Satisfies Ideal Bungee Physics]
  \label{def:physics:phyx-mini-0718:phyxminiproblems-problemphyxmini0718-satisfiesidealbungeephysics}
  \lean{PhyXMiniProblems.ProblemPhyXMini0718.SatisfiesIdealBungeePhysics}
  \uses{def:physics:phyx-mini-0718:phyxminiproblems-problemphyxmini0718-massinkilograms, def:physics:phyx-mini-0718:phyxminiproblems-problemphyxmini0718-lengthinmeters, def:physics:phyx-mini-0718:phyxminiproblems-problemphyxmini0718-timeinseconds, def:physics:phyx-mini-0718:phyxminiproblems-problemphyxmini0718-stiffnessinnewtonspermeter, def:physics:phyx-mini-0718:phyxminiproblems-problemphyxmini0718-accelerationinmeterspersecondsquared, def:physics:phyx-mini-0718:phyxminiproblems-problemphyxmini0718-signedlengthinmeters, def:physics:phyx-mini-0718:phyxminiproblems-problemphyxmini0718-velocityinmeterspersecond, def:physics:phyx-mini-0718:phyxminiproblems-problemphyxmini0718-bungeeoscillationsetup}
  The ideal vertical spring model consists of: static force balance k x\_eq = m g; amplitude equal to the excess of the release extension over the equilibrium extension; release from rest at the displayed point y = -A; Physlib's harmonic oscillator with SI-calibrated mass and stiffness; displacement and velocity given, for every elapsed duration, by Physlib's harmonic-oscillator trajectory and its time derivative. These are generic physical/modeling relations.
\end{definition}
\begin{proof}
  This is the declared model definition of Satisfies Ideal Bungee Physics.
\end{proof}
\begin{definition}[Answer Choice]
  \label{def:physics:phyx-mini-0718:phyxminiproblems-problemphyxmini0718-answerchoice}
  \lean{PhyXMiniProblems.ProblemPhyXMini0718.AnswerChoice}
  Letter labels of the four velocities printed with the problem
\end{definition}
\begin{proof}
  This is the declared model definition of Answer Choice.
\end{proof}
\begin{definition}[velocity Meters Per Second]
  \label{def:physics:phyx-mini-0718:phyxminiproblems-problemphyxmini0718-answerchoice-velocitymeterspersecond}
  \lean{PhyXMiniProblems.ProblemPhyXMini0718.AnswerChoice.velocityMetersPerSecond}
  \uses{def:physics:phyx-mini-0718:phyxminiproblems-problemphyxmini0718-answerchoice}
  Metre-per-second value printed beside each displayed answer label
\end{definition}
\begin{proof}
  This is the declared model definition of velocity Meters Per Second.
\end{proof}
\begin{definition}[recorded Answer Choice]
  \label{def:physics:phyx-mini-0718:phyxminiproblems-problemphyxmini0718-recordedanswerchoice}
  \lean{PhyXMiniProblems.ProblemPhyXMini0718.recordedAnswerChoice}
  \uses{def:physics:phyx-mini-0718:phyxminiproblems-problemphyxmini0718-answerchoice}
  The answer label recorded by the supplied dataset metadata
\end{definition}
\begin{proof}
  This is the declared model definition of recorded Answer Choice.
\end{proof}
% --- Archon declaration topology end ---
% --- Archon physics formalization source end ---
